% AdaBoost turns classification outcomes into the next round's sample weights.
% Schematic only; no empirical quantities are encoded.
% Canvas: 112 mm x 61 mm. Dependencies: project palette and TikZ.
\begin{tikzpicture}[
  x=1mm,y=1mm,
  font=\figurefont\footnotesize,
  text=FigureInk,
  stage/.style={
    rectangle,
    model hidden,
    minimum width=24mm,
    minimum height=11mm,
    text width=21mm,
    align=center,
    inner sep=1mm
  },
  data/.style={
    stage,
    model input
  },
  model/.style={
    stage,
    model hidden
  },
  update/.style={
    rectangle,
    model training,
    minimum width=29mm,
    minimum height=11mm,
    text width=26mm,
    align=center,
    inner sep=1mm
  },
  flow/.style={
    model flow
  },
  feedback/.style={
    model feedback
  }
]
  \path[use as bounding box] (0,0) rectangle (112,61);

  \node[anchor=west,font=\figurefont\small] at (1,57)
    {第$t$轮：分类结果决定第$t+1$轮的训练分布};

  \node[data] (weighted-data) at (15,40)
    {带权训练样本\\$\{w_i^{(t)}\}$};
  \node[model] (learner) at (43,40)
    {训练弱分类器\\$h_t$};
  \node[update] (outcome) at (88,49)
    {正确分类\\乘以$e^{-\alpha_t}$};
  \node[update] (mistake) at (88,31)
    {错分\\乘以$e^{\alpha_t}$};
  \node[update,minimum width=29mm] (normalize) at (44,12)
    {归一化\\得到$\{w_i^{(t+1)}\}$};

  \draw[flow] (weighted-data)--(learner);
  \draw[FigureInk,line width=.9pt] (learner.east)--(63,40);
  \fill[FigureInk] (63,40) circle[radius=1.1pt];
  \draw[flow] (63,40)|-(outcome.west);
  \draw[feedback] (63,40)|-(mistake.west);

  \draw[FigureInk,line width=.9pt]
    (outcome.east)--(109,49)--(109,20)--(106,20);
  \draw[FigureRed,dashed,line width=.9pt]
    (mistake.east)--(106,31)--(106,20);
  \fill[FigureInk] (106,20) circle[radius=1.1pt];
  \draw[flow] (106,20)--(65,20)--(65,12)--(normalize.east);
  \draw[feedback] (normalize.west)--(1,12)|-(weighted-data.west);
\end{tikzpicture}
